% Options for packages loaded elsewhere
\PassOptionsToPackage{unicode}{hyperref}
\PassOptionsToPackage{hyphens}{url}
\documentclass[
]{article}
\usepackage{xcolor}
\usepackage[margin=1in]{geometry}
\usepackage{amsmath,amssymb}
\setcounter{secnumdepth}{-\maxdimen} % remove section numbering
\usepackage{iftex}
\ifPDFTeX
  \usepackage[T1]{fontenc}
  \usepackage[utf8]{inputenc}
  \usepackage{textcomp} % provide euro and other symbols
\else % if luatex or xetex
  \usepackage{unicode-math} % this also loads fontspec
  \defaultfontfeatures{Scale=MatchLowercase}
  \defaultfontfeatures[\rmfamily]{Ligatures=TeX,Scale=1}
\fi
\usepackage{lmodern}
\ifPDFTeX\else
  % xetex/luatex font selection
\fi
% Use upquote if available, for straight quotes in verbatim environments
\IfFileExists{upquote.sty}{\usepackage{upquote}}{}
\IfFileExists{microtype.sty}{% use microtype if available
  \usepackage[]{microtype}
  \UseMicrotypeSet[protrusion]{basicmath} % disable protrusion for tt fonts
}{}
\makeatletter
\@ifundefined{KOMAClassName}{% if non-KOMA class
  \IfFileExists{parskip.sty}{%
    \usepackage{parskip}
  }{% else
    \setlength{\parindent}{0pt}
    \setlength{\parskip}{6pt plus 2pt minus 1pt}}
}{% if KOMA class
  \KOMAoptions{parskip=half}}
\makeatother
\usepackage{color}
\usepackage{fancyvrb}
\newcommand{\VerbBar}{|}
\newcommand{\VERB}{\Verb[commandchars=\\\{\}]}
\DefineVerbatimEnvironment{Highlighting}{Verbatim}{commandchars=\\\{\}}
% Add ',fontsize=\small' for more characters per line
\usepackage{framed}
\definecolor{shadecolor}{RGB}{248,248,248}
\newenvironment{Shaded}{\begin{snugshade}}{\end{snugshade}}
\newcommand{\AlertTok}[1]{\textcolor[rgb]{0.94,0.16,0.16}{#1}}
\newcommand{\AnnotationTok}[1]{\textcolor[rgb]{0.56,0.35,0.01}{\textbf{\textit{#1}}}}
\newcommand{\AttributeTok}[1]{\textcolor[rgb]{0.13,0.29,0.53}{#1}}
\newcommand{\BaseNTok}[1]{\textcolor[rgb]{0.00,0.00,0.81}{#1}}
\newcommand{\BuiltInTok}[1]{#1}
\newcommand{\CharTok}[1]{\textcolor[rgb]{0.31,0.60,0.02}{#1}}
\newcommand{\CommentTok}[1]{\textcolor[rgb]{0.56,0.35,0.01}{\textit{#1}}}
\newcommand{\CommentVarTok}[1]{\textcolor[rgb]{0.56,0.35,0.01}{\textbf{\textit{#1}}}}
\newcommand{\ConstantTok}[1]{\textcolor[rgb]{0.56,0.35,0.01}{#1}}
\newcommand{\ControlFlowTok}[1]{\textcolor[rgb]{0.13,0.29,0.53}{\textbf{#1}}}
\newcommand{\DataTypeTok}[1]{\textcolor[rgb]{0.13,0.29,0.53}{#1}}
\newcommand{\DecValTok}[1]{\textcolor[rgb]{0.00,0.00,0.81}{#1}}
\newcommand{\DocumentationTok}[1]{\textcolor[rgb]{0.56,0.35,0.01}{\textbf{\textit{#1}}}}
\newcommand{\ErrorTok}[1]{\textcolor[rgb]{0.64,0.00,0.00}{\textbf{#1}}}
\newcommand{\ExtensionTok}[1]{#1}
\newcommand{\FloatTok}[1]{\textcolor[rgb]{0.00,0.00,0.81}{#1}}
\newcommand{\FunctionTok}[1]{\textcolor[rgb]{0.13,0.29,0.53}{\textbf{#1}}}
\newcommand{\ImportTok}[1]{#1}
\newcommand{\InformationTok}[1]{\textcolor[rgb]{0.56,0.35,0.01}{\textbf{\textit{#1}}}}
\newcommand{\KeywordTok}[1]{\textcolor[rgb]{0.13,0.29,0.53}{\textbf{#1}}}
\newcommand{\NormalTok}[1]{#1}
\newcommand{\OperatorTok}[1]{\textcolor[rgb]{0.81,0.36,0.00}{\textbf{#1}}}
\newcommand{\OtherTok}[1]{\textcolor[rgb]{0.56,0.35,0.01}{#1}}
\newcommand{\PreprocessorTok}[1]{\textcolor[rgb]{0.56,0.35,0.01}{\textit{#1}}}
\newcommand{\RegionMarkerTok}[1]{#1}
\newcommand{\SpecialCharTok}[1]{\textcolor[rgb]{0.81,0.36,0.00}{\textbf{#1}}}
\newcommand{\SpecialStringTok}[1]{\textcolor[rgb]{0.31,0.60,0.02}{#1}}
\newcommand{\StringTok}[1]{\textcolor[rgb]{0.31,0.60,0.02}{#1}}
\newcommand{\VariableTok}[1]{\textcolor[rgb]{0.00,0.00,0.00}{#1}}
\newcommand{\VerbatimStringTok}[1]{\textcolor[rgb]{0.31,0.60,0.02}{#1}}
\newcommand{\WarningTok}[1]{\textcolor[rgb]{0.56,0.35,0.01}{\textbf{\textit{#1}}}}
\usepackage{graphicx}
\makeatletter
\newsavebox\pandoc@box
\newcommand*\pandocbounded[1]{% scales image to fit in text height/width
  \sbox\pandoc@box{#1}%
  \Gscale@div\@tempa{\textheight}{\dimexpr\ht\pandoc@box+\dp\pandoc@box\relax}%
  \Gscale@div\@tempb{\linewidth}{\wd\pandoc@box}%
  \ifdim\@tempb\p@<\@tempa\p@\let\@tempa\@tempb\fi% select the smaller of both
  \ifdim\@tempa\p@<\p@\scalebox{\@tempa}{\usebox\pandoc@box}%
  \else\usebox{\pandoc@box}%
  \fi%
}
% Set default figure placement to htbp
\def\fps@figure{htbp}
\makeatother
\setlength{\emergencystretch}{3em} % prevent overfull lines
\providecommand{\tightlist}{%
  \setlength{\itemsep}{0pt}\setlength{\parskip}{0pt}}
\usepackage{bookmark}
\IfFileExists{xurl.sty}{\usepackage{xurl}}{} % add URL line breaks if available
\urlstyle{same}
\hypersetup{
  pdftitle={Posterior Predictive Checks},
  hidelinks,
  pdfcreator={LaTeX via pandoc}}

\title{Posterior Predictive Checks}
\author{}
\date{\vspace{-2.5em}2026-04-17}

\begin{document}
\maketitle

\subsection{Introduction}\label{introduction}

A posterior predictive check is the Bayesian answer to ``does this model
produce data that look like ours?'' After fitting, you draw replicated
datasets from the model --- each one using a different posterior
parameter draw --- and compare them to the observed data through summary
statistics, density overlays, or grouped subsets. Systematic
discrepancies between observed and replicated quantities are a model
criticism: they tell you which feature of the data the current
specification fails to capture, which is far more actionable than a
single goodness-of-fit \(p\)-value.

\subsection{Prerequisites}\label{prerequisites}

A working understanding of the posterior predictive distribution and
access to MCMC draws or any object that exposes a
\texttt{posterior\_predict()} method (\texttt{brms}, \texttt{rstanarm},
raw Stan).

\subsection{Theory}\label{theory}

The posterior predictive distribution is

\[p(\tilde y \mid y) = \int p(\tilde y \mid \theta) p(\theta \mid y) \, \mathrm d \theta,\]

approximated by drawing \(\theta\) from the posterior and then
\(\tilde y\) from the likelihood given that \(\theta\). A posterior
predictive check chooses a test statistic \(T(y)\) --- the mean, the
standard deviation, the proportion of zeros, the maximum, the skew, or
any function of the data --- and compares its observed value to the
distribution of \(T(\tilde y)\) across replicates. If the observed
statistic falls in the bulk of the replicate distribution, the model
captures that aspect of the data; if it falls in a tail, the model
misses it.

\subsection{Assumptions}\label{assumptions}

A converged posterior with adequate effective sample size, since each
replicate dataset depends on a posterior draw. PPCs do not require
held-out data and are therefore not a measure of generalisation; they
are a within-sample diagnostic.

\subsection{R Implementation}\label{r-implementation}

\begin{Shaded}
\begin{Highlighting}[]
\FunctionTok{library}\NormalTok{(brms); }\FunctionTok{library}\NormalTok{(bayesplot)}

\FunctionTok{set.seed}\NormalTok{(}\DecValTok{2026}\NormalTok{)}
\NormalTok{d }\OtherTok{\textless{}{-}} \FunctionTok{data.frame}\NormalTok{(}\AttributeTok{y =} \FunctionTok{rnorm}\NormalTok{(}\DecValTok{200}\NormalTok{, }\DecValTok{3}\NormalTok{, }\FloatTok{1.5}\NormalTok{))}

\NormalTok{fit }\OtherTok{\textless{}{-}} \FunctionTok{brm}\NormalTok{(y }\SpecialCharTok{\textasciitilde{}} \DecValTok{1}\NormalTok{, }\AttributeTok{data =}\NormalTok{ d, }\AttributeTok{family =}\NormalTok{ gaussian,}
           \AttributeTok{chains =} \DecValTok{2}\NormalTok{, }\AttributeTok{iter =} \DecValTok{2000}\NormalTok{, }\AttributeTok{refresh =} \DecValTok{0}\NormalTok{)}

\CommentTok{\# Density overlay}
\FunctionTok{pp\_check}\NormalTok{(fit)}

\CommentTok{\# Test{-}statistic check: SD of simulated vs observed}
\FunctionTok{pp\_check}\NormalTok{(fit, }\AttributeTok{type =} \StringTok{"stat"}\NormalTok{, }\AttributeTok{stat =} \StringTok{"sd"}\NormalTok{)}

\CommentTok{\# Interval check}
\FunctionTok{pp\_check}\NormalTok{(fit, }\AttributeTok{type =} \StringTok{"intervals"}\NormalTok{)}
\end{Highlighting}
\end{Shaded}

\subsection{Output \& Results}\label{output-results}

\texttt{pp\_check(fit)} returns a \texttt{ggplot} overlaying the
empirical density of \texttt{y} on top of densities for several
replicate datasets. \texttt{type\ =\ "stat"} plots the distribution of
the chosen statistic across replicates with the observed value as a
vertical line. \texttt{type\ =\ "intervals"} shows the 50 \% and 90 \%
posterior predictive intervals point-by-point and is the most
informative check for time-series and ordered data.

\subsection{Interpretation}\label{interpretation}

A reporting sentence: ``Posterior predictive density overlays showed the
observed distribution falling well within the cloud of replicates; the
standard-deviation check placed the observed value near the centre of
the replicate distribution (\(p_{\text{ppc}} = 0.48\)), indicating that
the Gaussian model captures both location and dispersion adequately.''
When a check fails, name the feature: ``skew of replicates was
systematically lower than the observed skew'', which directly suggests a
heavy-tailed family or a transformation.

\subsection{Practical Tips}\label{practical-tips}

\begin{itemize}
\tightlist
\item
  A density overlay is the broadest check; supplement it with tests on
  specific statistics (skew, kurtosis, proportion of zeros) that target
  likely misspecifications.
\item
  Group-wise PPCs
  (\texttt{pp\_check(fit,\ type\ =\ "stat\_grouped",\ group\ =\ "site")})
  catch problems concentrated in specific subgroups that pooled checks
  would mask.
\item
  For count outcomes, \texttt{type\ =\ "rootogram"} is more informative
  than density overlays because it highlights overdispersion and
  zero-inflation directly.
\item
  A passing PPC does not guarantee good predictions on new data; pair
  PPCs with \texttt{loo()} or held-out evaluation when generalisation
  matters.
\item
  When PPCs reveal misspecification, the fix is almost always a richer
  family (heavy tails, dispersion, mixtures) before a richer mean
  structure.
\end{itemize}

\subsection{R Packages Used}\label{r-packages-used}

\texttt{bayesplot} for the \texttt{pp\_check()} and \texttt{ppc\_*}
graphics, \texttt{brms} (or \texttt{rstanarm}) as a backend that exposes
\texttt{posterior\_predict()} directly, and \texttt{posterior} for
handling draw arrays when building custom PPCs by hand.

\subsection{For Reviewers}\label{for-reviewers}

\textbf{What to look for in a paper using this method.}

\begin{itemize}
\tightlist
\item
  \textbf{Common misapplications.}
\item
  \textbf{Diagnostics that should be reported but often aren't.}
\item
  \textbf{Red flags in tables and figures.}
\item
  \textbf{What to verify.}
\item
  \textbf{What an adequate Methods paragraph must contain.}
\end{itemize}

\subsection{See also --- labs in this
chapter}\label{see-also-labs-in-this-chapter}

\begin{itemize}
\tightlist
\item
  \href{labs/lab-bayesian-thinking-control-vs-decision-error.Rmd}{Bayesian
  thinking; α control vs decision error}
\item
  \href{labs/lab-brms-stan-loo-hierarchical-models.Rmd}{brms / Stan,
  LOO, hierarchical models}
\end{itemize}

\end{document}
